\section{Before-After Comparison}

Table~\ref{tab:before-after} presents the corridor scores under v1.0 (120-point scale), v1.1 (120-point, A3 corrected), and v1.2 (150-point, strategic dimensions added), along with the tier assignment at each version.

\begin{longtable}{lrrlrrlrrl}
\caption{Corridor Scores and Tiers Across Rubric Versions} \label{tab:before-after} \\
\toprule
\textbf{Corridor} & \textbf{v1.0} & \textbf{T} & \textbf{v1.1} & \textbf{T} & \textbf{v1.2} & \textbf{T} & \textbf{Key change} \\
\midrule
\endfirsthead
\multicolumn{8}{c}{Table~\ref{tab:before-after} continued} \\
\toprule
\textbf{Corridor} & \textbf{v1.0} & \textbf{T} & \textbf{v1.1} & \textbf{T} & \textbf{v1.2} & \textbf{T} & \textbf{Key change} \\
\midrule
\endhead
\bottomrule
\endfoot

\multicolumn{8}{l}{\textit{Corrected T1 arteries (gained from v1.0)}} \\
I-80 & 25.2 & T2 & 20.2 & T2 & 38.7 & \textbf{T1} & A4+B4+C4 add 18.5pts \\
I-10 & 26.5 & T1 & 21.5 & T1 & 37.9 & \textbf{T1} & All three dimensions contribute \\
I-5  & 26.8 & T1 & 21.8 & T1 & 36.0 & \textbf{T1} & USMCA San Diego highest \\
I-35 & -- & T3 & 12.2 & T3 & 33.8 & \textbf{T1} & Laredo A4=8.5 + grain C4=8.0 \\
I-90 & 28.0 & T1 & 23.0 & T1 & 32.8 & \textbf{T1} & ICBM B4=9.0 decisive \\
I-95 & 24.8 & T2 & 19.8 & T2 & 21.2 & T2* & No ag/trade; strategic override \\
I-75 & 18.8 & T2 & 13.8 & T3 & 23.7 & T2* & strategic override \\
\midrule
\multicolumn{8}{l}{\textit{Demoted (were T1, now correctly lower)}} \\
I-110 & 30.1 & T1 & 25.1 & T1 & 19.0 & T2 & A3 fix removes 5pts; no strategic dims \\
I-880 & 30.0 & T1 & 25.0 & T1 & 15.0 & T4 & Same; drops to T4 \\
I-225 & 29.6 & T1 & 24.6 & T1 & 14.6 & T3 & Denver beltway; no strategic value \\
\midrule
\multicolumn{8}{l}{\textit{US highway upgrade candidates (new in v1.2)}} \\
US-287 & -- & -- & 14.0 & T3 & 35.5 & T1-range & B1=10 + B4=7.5 + C4=9.0 \\
US-2   & -- & -- & 33.5 & -- & 33.5 & T1-range & B1=10; TIGER only 31mi \\
US-83  & -- & -- & 27.0 & -- & 27.0 & T1-range & USMCA Eagle Pass A4=7.5 \\
\end{longtable}

*T1 by strategic override; v1.2 score below threshold due to C4=0 (correct data for urban corridors) but national strategic role is unambiguous.

\subsection{The Congestion-Stress Paradox: Quantified}

The clearest demonstration of the v1.2 correction is the I-110/I-80 reversal:

\begin{itemize}
\item v1.0: I-110 = 30.1/120 (T1 \#1), I-80 = 25.2/120 (T2)
\item v1.2: I-80 = 38.7/150 (T1 \#1), I-110 = 19.0/150 (T2)
\end{itemize}

The reversal occurred because:
\begin{enumerate}
\item A3 fix: I-110's IRI-driven A3=10.0 dropped to 5.0; I-80's dropped the same amount (net neutral between them)
\item A4+B4+C4 and recalibration: I-80 gained 18.5 points total --- A4=0.0 (no primary border crossing), B4=6.5 (Cheyenne/FE Warren proximity), C4=7.5 (Corn Belt, dual-coast export access), with the remaining 4.5 points from A3 BPR path improvement and B2 recalibration. I-110 gained 5.0 from STRAHNET baseline only (no trade, no ag)
\end{enumerate}

The 19.7-point gap between I-80 and I-110 in v1.2 (vs. 4.9 points in v1.0 going the other way) reflects the actual strategic difference between a transcontinental trunk line and an urban connector.

\subsection{US-287: The Most Dramatic Gain}

US-287 went from 14.0/120 in v1.1 (where it was invisible as a T3 corridor) to 35.5/150 in v1.2 --- within 3 points of I-80. This occurred because:

\begin{itemize}
\item B1=10.0: Maximum isolation; no parallel interstate within 200 miles across most of its alignment
\item B4=7.5: FE Warren AFB (USSTRATCOM approach); Fort Carson CO
\item C4=9.0: Texas Panhandle sorghum (\#1 US producer); Wyoming cattle; dual-access to Gulf AND Pacific ports
\end{itemize}

The v1.1 rubric was blind to all of this. US-287 carried the same strategic value in v1.1 that it carries in v1.2 --- the rubric simply lacked the dimensions to measure it.

\subsection{What Didn't Change}

Fourteen corridors maintained consistent tier assignments across all three rubric versions. These corridors have unambiguous status regardless of rubric calibration:
\begin{itemize}
\item Stable T1: I-90, I-5, I-10 (all three versions)
\item Stable T3/T4: Short urban spurs (I-184, I-194, etc.) — consistently low regardless of strategic dimensions
\end{itemize}

Stability under different rubric versions is evidence of correct classification. The corridors that \textit{changed} tier across versions are the ones the rubric most needed to correct.
